\documentclass[10pt]{article}
% Shared LaTeX style for MIT 18.06 exam revision notes.
% Compile topic files with Tectonic, XeLaTeX, or LuaLaTeX.

\usepackage[a4paper,margin=0.72in]{geometry}
\usepackage{fontspec}
\usepackage{amsmath,amssymb,mathtools}
\usepackage{booktabs,array,tabularx}
\usepackage{enumitem}
\usepackage{titlesec}
\usepackage{fancyhdr}
\usepackage{microtype}
\usepackage{xcolor}

\setmainfont{Times New Roman}
\setsansfont{Arial}

\setlength{\parindent}{0pt}
\setlength{\parskip}{3.6pt}
\setlength{\abovedisplayskip}{6pt}
\setlength{\belowdisplayskip}{6pt}
\setlength{\abovedisplayshortskip}{4pt}
\setlength{\belowdisplayshortskip}{4pt}
\renewcommand{\arraystretch}{1.08}
\setlength{\tabcolsep}{4.5pt}

\titleformat{\section}
  {\normalfont\large\bfseries}
  {\thesection}
  {0.55em}
  {}
\titlespacing*{\section}{0pt}{10pt}{3pt}

\titleformat{\subsection}
  {\normalfont\normalsize\bfseries}
  {\thesubsection}
  {0.5em}
  {}
\titlespacing*{\subsection}{0pt}{7pt}{2pt}

\setlist[itemize]{leftmargin=1.35em,itemsep=1pt,topsep=2pt,parsep=0pt}
\setlist[enumerate]{leftmargin=1.55em,itemsep=1pt,topsep=2pt,parsep=0pt}

\pagestyle{fancy}
\fancyhf{}
\fancyfoot[C]{\scriptsize\color{black!45}MIT 18.06 Linear Algebra Exam Notes --- \thepage}
\renewcommand{\headrulewidth}{0pt}
\renewcommand{\footrulewidth}{0pt}

\newcommand{\notesubtitle}[1]{%
  {\centering\itshape #1\par}
  \vspace{2pt}
}

\newcommand{\source}[1]{%
  {\centering\small #1\par}
  \vspace{6pt}
}

\newcommand{\labelnote}[2]{%
  \par\smallskip
  \noindent\textbf{#1.}\quad #2
  \par\smallskip
}

\newcommand{\tightlabel}[1]{%
  \par\smallskip
  \noindent\textbf{#1.}\quad
}

\newcolumntype{Y}{>{\raggedright\arraybackslash}X}
\newcolumntype{L}[1]{>{\raggedright\arraybackslash}p{#1}}

\newenvironment{notetable}
  {\small\begin{tabularx}{\linewidth}}
  {\end{tabularx}}


\begin{document}

\begin{center}
{\LARGE 4.4, 5.1--5.2 Orthonormal Bases, QR, and Determinants\par}
\end{center}
\notesubtitle{Concise exam revision notes for orthonormal columns, QR factorization, and determinant rules.}
\source{Based on handwritten MIT 18.06 revision notes.}

\section{Core Idea}

\labelnote{Core idea}{Orthonormal columns make projection and least squares simple. Determinants measure whether a square matrix is singular and how elimination changes volume.}

This note has two connected parts:
\[
Q^{T}Q=I
\qquad\text{and}\qquad
\det A.
\]
The first simplifies geometry; the second summarizes invertibility for square matrices.

\section{Orthogonal vs Orthonormal}

\begin{center}
\small
\begin{tabularx}{\linewidth}{@{}L{0.25\linewidth}Y Y@{}}
\toprule
Word & Condition & Meaning \\
\midrule
Orthogonal & \(q_i^{T}q_j=0\) for \(i\ne j\). & Different directions are perpendicular. \\
Orthonormal & \(q_i^{T}q_j=0\) and \(q_i^{T}q_i=1\). & Perpendicular and unit length. \\
Matrix \(Q\) with orthonormal columns & \(Q^{T}Q=I\). & Dot products are already built in. \\
\bottomrule
\end{tabularx}
\end{center}

If \(Q\) is square and orthonormal, then
\[
Q^{-1}=Q^{T}.
\]
For rectangular \(Q\), \(Q^{T}Q=I\), but \(QQ^{T}\) is a projection matrix.

\section{Gram--Schmidt}

Gram--Schmidt turns independent vectors into orthogonal vectors, then normalizes them.

For two vectors \(a,b\), first keep
\[
A=a.
\]
Remove from \(b\) the part in the \(a\)-direction:
\[
B=b-\frac{a^{T}b}{a^{T}a}a.
\]
Then normalize:
\[
q_1=\frac{A}{\|A\|},
\qquad
q_2=\frac{B}{\|B\|}.
\]
For a third vector \(c\), subtract the projections onto the previous orthogonal directions:
\[
C=c-\frac{A^{T}c}{A^{T}A}A-\frac{B^{T}c}{B^{T}B}B.
\]

\labelnote{Exam trigger}{If the problem says ``make an orthonormal basis'', use Gram--Schmidt, then divide by lengths.}

\section{QR Factorization}

For independent columns,
\[
A=QR,
\]
where \(Q\) has orthonormal columns and \(R\) is upper triangular.

The entries of \(R\) record dot products with the \(q\)'s:
\[
R=Q^{T}A.
\]
For least squares,
\[
A^{T}A\hat{x}=A^{T}b
\]
becomes
\[
R\hat{x}=Q^{T}b.
\]
This is easier because \(R\) is triangular.

\section{Projection with Orthonormal Columns}

If \(Q\) has orthonormal columns, projection onto \(C(Q)\) is
\[
p=QQ^{T}b.
\]
In vector form,
\[
p=(q_1^{T}b)q_1+(q_2^{T}b)q_2+\cdots+(q_n^{T}b)q_n.
\]

\begin{center}
\small
\begin{tabularx}{\linewidth}{@{}L{0.28\linewidth}Y@{}}
\toprule
Column type & Projection habit \\
\midrule
Orthonormal \(q_i\) & Use \(q_i^{T}b\) directly. \\
Orthogonal but not unit & Divide by \(a_i^{T}a_i\). \\
General independent columns & Use \(A(A^{T}A)^{-1}A^{T}\). \\
\bottomrule
\end{tabularx}
\end{center}

\section{Orthogonal Matrices}

For square orthogonal \(Q\),
\[
Q^{T}Q=I,
\qquad
Q^{-1}=Q^{T}.
\]
Orthogonal matrices preserve lengths and dot products:
\[
\|Qx\|=\|x\|,
\qquad
(Qx)^{T}(Qy)=x^{T}y.
\]
Rotations and reflections are the main exam examples.

\section{Determinant Rules}

\begin{center}
\small
\begin{tabularx}{\linewidth}{@{}L{0.34\linewidth}Y@{}}
\toprule
Operation or fact & Effect on determinant \\
\midrule
Exchange two rows & Changes the sign. \\
Multiply one row by \(c\) & Multiplies determinant by \(c\). \\
Add a multiple of one row to another & No change. \\
Triangular matrix & Product of diagonal entries. \\
Singular matrix & Determinant is \(0\). \\
Product & \(\det(AB)=\det A\det B\). \\
Transpose & \(\det(A^{T})=\det A\). \\
\bottomrule
\end{tabularx}
\end{center}

\labelnote{Exam memory}{A zero determinant means no inverse. A nonzero determinant means the square matrix is invertible.}

\section{Pivots and Cofactors}

If elimination uses no row exchanges, then
\[
\det A=(\text{product of pivots}).
\]
For leading submatrices \(A_k\), pivots can be checked by determinant ratios:
\[
d_k=\frac{\det A_k}{\det A_{k-1}}.
\]

Cofactor expansion is useful when a row or column has many zeros:
\[
\det A=\sum_j (-1)^{i+j}a_{ij}\det M_{ij}.
\]
Use it selectively; elimination is usually faster for larger matrices.

\section{Common Mistakes}

\begin{center}
\small
\begin{tabularx}{\linewidth}{@{}L{0.40\linewidth}Y@{}}
\toprule
Mistake & Correct exam habit \\
\midrule
Calling orthogonal columns orthonormal. & Check both perpendicularity and unit length. \\
Using \(Q^{-1}=Q^{T}\) for rectangular \(Q\). & Only square orthogonal matrices have inverse \(Q^{T}\). \\
Forgetting to normalize after Gram--Schmidt. & Orthogonal first, unit length second. \\
Using \(QQ^{T}\) for any basis. & Use it only for orthonormal columns. \\
Changing a row and forgetting the determinant effect. & Track swaps, scaling, and row replacement separately. \\
Expanding cofactors when elimination is simpler. & Use zeros or structure to choose the method. \\
\bottomrule
\end{tabularx}
\end{center}

\section{Final Exam Checklist}

\tightlabel{Checklist}
\begin{enumerate}
  \item Check whether vectors are orthogonal by dot products.
  \item Check whether vectors are orthonormal by dot products and lengths.
  \item Use Gram--Schmidt to build orthogonal vectors.
  \item Normalize each vector to get \(q_i\).
  \item Form \(Q\) from the \(q_i\)'s and compute \(R=Q^{T}A\).
  \item For least squares with \(A=QR\), solve \(R\hat{x}=Q^{T}b\).
  \item For determinants, choose elimination unless cofactor expansion has many zeros.
  \item Track determinant changes from row operations.
  \item Use \(\det A=0\) to identify singular square matrices.
\end{enumerate}

\end{document}
